% Part: incompleteness
% Chapter: theories-computability
% Section: oconsis-ext-of-q-undec

\documentclass[../../../include/open-logic-section]{subfiles}

\begin{document}

\olfileid{inc}{tcp}{oqn}

\olsection{Yen $\omega$-konsisten, ekstensi $\Th{Q}$ ora bisa diputusake}

\begin{explain}
Bukti yen $\Th{Q}$ iku c.e.-lengkap gumantung marang kasunyatan yen saben
ukara kang bisa dibuktekake ing $\Th{Q}$ iku ``bener'' tumrap bilangan asli.
Definisi lan teorema sabanjure ngiyatake teorema iki kanthi nemtokake
persis bagean saka ``bebener'' kang dibutuhake ing bukti ing
ndhuwur. Nanging aja suwe-suwe mandheg ing teorema iki, amarga ora
let suwe maneh kita bakal ngiyatake luwih adoh. Iki kalebu utamane
amarga alasan sejarah: makalah asli G\"odel nggunakake gagasan
kekonsistenan-$\omega$, nanging asile enggal dikuatake kanthi ngganti
kekonsistenan-$\omega$ nganggo kekonsistenan lumrah.
\end{explain}

\begin{defn}
\ollabel{thm:oconsis-q}
Teori $\Th{T}$ iku $\omega$-konsisten yen kahanan iki kelakon: yen
$\lexists[x][!A(x)]$ iku ukara sembarang lan $\Th{T}$ mbuktekake $\lnot
!A(\num 0)$, $\lnot !A(\num 1)$, $\lnot !A(\num 2)$, \dots mula $\Th{T}$
ora mbuktekake $\lexists[x][!A(x)]$.
\end{defn}

\begin{thm}
  Upamana $\Th{T}$ teori $\omega$-konsisten sembarang kang ngemot $\Th{Q}$.
  Mula $\Th{T}$ ora bisa diputusake.
\end{thm}

\begin{proof}
Yen $\Th{T}$ ngemot $\Th{Q}$, mula $\Th{T}$ merepresentasikake fungsi
lan relasi kang komputabel. Kita mung perlu ngowahi bukti sadurunge.
Kaya ing ndhuwur, yen $x \in K$, mula $\Th{T}$ mbuktekake $\lexists[s][!A_T(\num
  x,\num x, s)]$. Kosok baline, upamana $\Th{T}$ mbuktekake
$\lexists[s][!A_T(\num x, \num x, s)]$. Mula $x$ mesthi kalebu ing $K$:
yen ora, ora ana komputasi mesin $x$ kang mandheg nalika nampa masukan
$x$; amarga $!A_T$ merepresentasikake relasi $T$ saka Kleene, $\Th{T}$
mbuktekake $\lnot !A_T(\num x, \num x, \num 0)$, $\lnot !A_T(\num x, \num
x, \num 1)$, \dots, saengga $\Th{T}$ dadi ora $\omega$-konsisten.
\end{proof}

\end{document}
